\documentclass{article}
\usepackage{geometry}
\geometry{a4paper,scale=0.8}
\usepackage[utf8]{inputenc}

\usepackage{amsmath, amssymb}
\usepackage{amsthm}
\usepackage{enumitem}

\newtheoremstyle{breakstyle}
    {5pt}
    {10pt}
    % {\normalfont}
    {\itshape}
    {0pt}
    {\bfseries}
    {\newline\indent}
    {0pt}
    {\thmname{#1}\thmnumber{ #2}\thmnote{ \textbf{[#3]}}}
\theoremstyle{breakstyle}
\newtheorem{exercise}{Exercise}[section]
\newtheorem{remark}{Remark}[section]

\newenvironment{solution}
  {\begin{proof}[Solution.]\renewcommand{\qedsymbol}{}}
  {\end{proof}}

% \let\ker\relax
% \DeclareMathOperator{\ker}{Ker}
\let\det\relax
\DeclareMathOperator{\det}{Det}

\title{Chapter 2 - Classification of finite reflection groups}
\author{Flandre}
\date{\today}

\begin{document}
\maketitle

\setcounter{section}{1}
\section{Irreducible components}

\begin{exercise}[Skipped]
Let $W$ be the dihedral group $\mathcal{D}_6$ of order 12, with standard Coxeter
generators $S = \{s, s'\}$ . The Coxeter system $(W, S)$ is irreducible.
However, $W$ has another set $S'$ of Coxeter generators leading to a Coxeter
system which is not irreducible: $S' := \{s, (s's)^3, s(s's)^2\}$ .
\end{exercise}

\begin{remark}
    This tells us that irreducibility is a property of the Coxeter system, i.e.
    of the chosen set $S$, not of the abstract group $W$.
\end{remark}

\setcounter{section}{7}
\section{Crystallographic groups}

\begin{exercise}
If $W$ satisfies the necessary condition in the proposition (and is essential),
show how to modify the lengths of the roots in a simple system $\Delta$ so that
the resulting $\mathbf{Z}$ -span is a lattice in $V$ stable under $W$ . [Use the
fact that the Coxeter graph has no circuits.]

\begin{solution}
    For any two $\alpha_{i},\alpha_{j}\in \Delta$, we'll modify their length
    so that $2\frac{\left( \alpha_{i},\alpha_{j} \right) }{\left(
    \alpha_{i},\alpha_{i} \right) }\in \mathbb{Z}$, in which case
    $$
    s_{\alpha_{i}}\left( \alpha_{j} \right) =a_{j}-2\frac{\left(
    \alpha_{i},\alpha_{j} \right) }{\left( \alpha_{i},\alpha_{i} \right)
    }\alpha_{j}
    $$
    lies in the $\mathbb{Z}$-span of $\alpha_{i},\alpha_{j}$.

    Notice
    $$
    2\frac{\left( \alpha_{i},\alpha_{j} \right) }{\left(
    \alpha_{i},\alpha_{i} \right) }=\frac{\left| \alpha_{j} \right| }{\left|
    \alpha_{i} \right| }\cdot 2\cos\left< \alpha_{i},\alpha_{j} \right> 
    $$
    while
    $$
    2\frac{\left( \alpha_{i},\alpha_{j} \right) }{\left(
    \alpha_{j},\alpha_{j} \right) }=\frac{\left| \alpha_{i} \right| }{\left|
    \alpha_{j} \right| }\cdot 2\cos\left< \alpha_{i},\alpha_{j} \right> 
    $$
    
    And since $m\left( \alpha,\beta \right) \in \left\{ 2,3,4,6 \right\} $ since
    $W$ is crystallographic,
    $$
    \cos\left< \alpha_{i},\alpha_{j} \right> =\cos{\frac{\pi}{m\left(
    \alpha_{i},\alpha_{j} \right) }}\in \left\{ 0,\frac{1}{2},\frac{\sqrt{2}
    }{2},\frac{\sqrt{3} }{2} \right\}\text{.   ($0$ is omitted as $\alpha_{i}\neq
    \alpha_{j}$)}
    $$
    We can have $2\cos\left< \alpha_{i},\alpha_{j} \right>
    =\sqrt{n_{ij}},n_{ij}\in \mathbb{Z} $. Now, we may let $\frac{\left| a_{i}
    \right| }{\left| a_{j} \right| }=\sqrt{n_{ij}} $, and therefore
    $$
    2\frac{\left( \alpha_{i},\alpha_{j} \right) }{\left( \alpha_{i},\alpha_{i}
    \right) }=\frac{1}{\sqrt{n_{ij}} }\sqrt{n_{ij}} =1\in \mathbb{Z}
    $$
    $$
    2\frac{\left( \alpha_{j},\alpha_{i} \right) }{\left( \alpha_{j},\alpha_{j}
    \right) }=\sqrt{n_{ij}} \sqrt{n_{ij}} =n_{ij}\in \mathbb{Z}
    .$$

    \emph{(Following from the hint)}

    Now, we may pick out any $\alpha_{i}\in \Delta$ as the starting vertex,
    unify it to lenght $1$, then scale all the vertices, denote one of them by
    $\alpha_{j}$, in its neighbourhood to the length satisfying $\frac{\left|
    \alpha_{i} \right| }{\left| \alpha_{j}\right| }=\sqrt{n_{ij}} $. Since
    there's no loop in the corresponding Coxeter Graph, such process can go
    through all the vertices in this finite Coxeter Graph without contradiction,
    and the $\mathbb{Z}$-span was obtained.
\end{solution}

\end{exercise}

\end{document}
